\documentclass[12pt]{article}
\usepackage[T1]{fontenc}
\usepackage[utf8]{inputenc}
\usepackage{lmodern}
\usepackage{textcomp}
\usepackage{enumitem}
\usepackage{geometry}
\geometry{margin=1in}
\begin{document}
\begin{center}
Answer Key - Computer Science - Graduate Level Assessment 13 \
\small (Answers are ordered exactly as in the question paper.)
\end{center}

\section*{Multiple Choice}
\begin{enumerate}[label=\arabic*.]
\item \textbf{Answer:} B
\\ \textit{Options:} A) Values of local variables; B) A heap area; C) The return address; D) Stack pointer for the calling activation record
\\ \textit{Note:} A heap area is not represented in a subroutine's activation record frame because the activation record is specifically designed to manage local variables, return addresses, and control information for a single invocation of a subroutine, while the heap is a separate memory area used for dynamic memory allocation.
\item \textbf{Answer:} D
\\ \textit{Options:} A) Restrict what operations/data the user can access; B) Determine if the user is an attacker; C) Flag the user if he/she misbehaves; D) Determine who the user is
\\ \textit{Note:} Authentication aims to determine who the user is by verifying their identity through credentials such as passwords, biometrics, or tokens, thereby ensuring that access to systems and data is granted only to authorized individuals.
\item \textbf{Answer:} A
\\ \textit{Options:} A) Recursive procedures; B) Arbitrary goto's; C) Two-dimensional arrays; D) Integer-valued functions
\\ \textit{Note:} Recursive procedures require stack-based storage allocation because each recursive call needs to maintain its own execution context, including local variables and return addresses, which is efficiently managed by the stack.
\item \textbf{Answer:} A
\\ \textit{Options:} A) xyz; B) xy; C) xxzy; D) xxxxy
\\ \textit{Note:} The string "xyz" cannot be generated by the given grammar because it requires the presence of an E production to derive the 'z', and the grammar ultimately leads to productions that cannot yield this exact sequence without additional characters or structures.
\item \textbf{Answer:} A
\\ \textit{Options:} A) By overwriting the return address to point to the location of that code; B) By writing directly to the instruction pointer register the address of the code; C) By writing directly to \%eax the address of the code; D) By changing the name of the running executable, stored on the stack
\\ \textit{Note:} A buffer overflow on the stack allows an attacker to overwrite the return address of a function with the address of their injected code, enabling the program to execute malicious instructions instead of returning to the intended location.
\end{enumerate}

\section*{True / False}
\begin{enumerate}[label=\arabic*.]
\setcounter{enumi}{5}
\item \textbf{Answer:} False
\\ \textit{Options:} A) True; B) False
\\ \textit{Note:} WPA primarily utilizes the Temporal Key Integrity Protocol (TKIP) for encryption and the Extensible Authentication Protocol (EAP) for authentication, rather than relying on the Lightweight EAP (LEAP) protocol.
\item \textbf{Answer:} True
\\ \textit{Options:} A) True; B) False
\\ \textit{Note:} To verify a digital signature, the recipient must use the sender's public key because it is mathematically linked to the private key used to create the signature, ensuring that the signature is authentic and has not been altered.
\item \textbf{Answer:} False
\\ \textit{Options:} A) True; B) False
\\ \textit{Note:} A proxy firewall operates at the application layer of the OSI model, not the physical layer, as it inspects and manages data packets at the application level rather than handling hardware-level data transmission.
\item \textbf{Answer:} False
\\ \textit{Options:} A) True; B) False
\\ \textit{Note:} WPA 2 is considered less secure than its successor, WPA 3, which includes enhanced security features and protections against modern threats.
\item \textbf{Answer:} True
\\ \textit{Options:} A) True; B) False
\\ \textit{Note:} Availability is often viewed as a functional requirement rather than a classic security property, which traditionally emphasizes confidentiality and integrity, thereby leading to the classification of the statement as true.
\end{enumerate}

\section*{Long Form Response}
\begin{enumerate}[label=\arabic*.]
\setcounter{enumi}{10}
\item \textbf{Answer:} def drop\_empty(dict 1): | return dict 1
\\ \textit{Note:} The provided function simply returns the original dictionary without modifying it, thus failing to implement the requirement of dropping empty items, which indicates a misunderstanding of the function's purpose.
\item \textbf{Answer:} def split\_upperstring(text): | return (re.findall('[A-Z][^A-Z]*', text))
\\ \textit{Note:} correct because it utilizes a regular expression to identify uppercase letters and split the string accordingly, ensuring that each segment starts with an uppercase letter followed by any number of non-uppercase characters.
\end{enumerate}

\vfill
\noindent\textit{End of Answer Key}
\end{document}